% AI-written, AI cross-reviewed; no human sign-off; not a claim that RH is proved
% Atlas unit P024. Local manuscript v3, item 7.1, heading on printed PDF p. 72.
% UpperTheory.V.Z.Win.Dn.ends0_alt
% paper2/inputs_one.tex:12–44. Verbatim excerpt; SHA-256 6c61e297636e859c1a31ee019e898f04c15b46b29037df23193db8c5df6f31ba.
% Original TeX uses the macros and cross-reference keys of the source manuscript.


% Start of atlas edition annotation. The following block is outside the original excerpt.
\paragraph{Status of named inputs in this reading.}
\noindent\nolinkurl{UpperTheory.V.Z.Win.Dn.ends0_alt} --- in-repo constant (\path{Development/AutomorphicArithmetic/AutomorphicVarieties/ModularityArgument/ThresholdArguments/DenseThresholdCases/DnEndsTop.lean:15--16}); not compared.\par
% End of atlas edition annotation.

\begin{formalstatement}[Two conclusions with a common mark]\label{pn:P024}
Assume the automorphic hypotheses \(H'\) together with weak descent,
cyclic base change and quadratic descent, the specified quaternion
existence theorem, and real semisimplicity of the cusp data.
\begin{enumerate}
\item For every admitted CM field \(K\), coefficient ring, nonzero level,
weight parameter, character and marked datum \(D\), every
\(\iota:\overline{\mathbb Q}_\ell\simeq\mathbb C\), and every valuation
above \(\ell\), there is a single \(g_1\) satisfying the marking equations
for \(D\), the finiteness condition \(\mathrm{FinEnd}\), and both automorphic
end conditions, of ranks three and two respectively.
\item Suppose also that \([K:\mathbb Q]=2\) and \(K\) contains a root
of \(X^2+1\). For every admitted host datum and host frame \(F\) there
is a marked \(g_1\) such that for each universal representation \(u\)
at that mark, its weight-zero part is finite as a module over the
coefficient ring. At every Type~I specialization \(r\) of \(u\), the
indicated matrix representation satisfies the prescribed
automorphy condition over \(K\) with twist \(-1\).
\end{enumerate}
\end{formalstatement}\nodestatus{P024}

The two-line declaration applies the linked-unramified result to the rank-two
seed and the dense rank-three theorem. Blueprint Z1.1.07~\bploc{r8:small:Z1.1.07} constructs
one level on which the sign kernel, the cyclotomic character and the
required character bounds coexist. The missing step in the written argument
is compatibility of the flags and ends for the same seeds followed by
their joint descent. Barnet-Lamb, Gee, Geraghty and Taylor,
\emph{Potential automorphy and change of weight}, Theorem~2.3.1,
lift an individual polarized representation under residual adequacy and
local component hypotheses. Theorem~2.4.2 addresses finiteness for a
specified global deformation problem. The common mark in two ranks and
the weight-zero ring require a further comparison of deformation problems.
\cite{r21-blggt}
